% Part: lambda-calculus
% Chapter: lambda-definablity
% Section: arithmetical-functions

\documentclass[../../../include/open-logic-section]{subfiles}

\begin{document}

\olfileid{lam}{rep}{arf}
\olsection{$\lambda$-可定义的算术函数}

\begin{prop}
  \ollabel{prop:succ-ld}
  后继函数~$\Succ$ 是 $\lambd$-可定义的。
\end{prop}

\begin{proof}
一个 $\lambd$-定义了后继函数的项为
\begin{align*}
  \fn{Succ} & \ident \lambd[a][\lambd[fx][f ({a} f x)]].
  \intertext{根据我们的约定，它是以下项的简写：}
  \fn{Succ} & \ident \lambd[a][\lambd[f][\lambd[x][(f (({a} f) x))]]].
  \intertext{$\fn{Succ}$ 是一个以数字~$a$ 为自变元的函数，它求值为另一个函数，即 $\lambd[fx][f ({a}
      f x)]$。该函数本身并不是丘奇数。然而，如果自变元 $a$ 是丘奇数，它就归约为一个丘奇数。考虑：}
   (\lambd[a][\lambd[fx][f ({a} f x)]])\,\num{n} & \redone
   \lambd[fx][f ({\num{n}} f x)].
   \intertext{内嵌的项 $\num{n}fx$ 是一个可归约式，因为 $\num{n}$ 就是 $\lambd[fx][f^nx]$。所以 $\num{n}fx \redone f^nx$，于是，对整个项我们有}
   \fn{Succ}\, \num n & \red \lambd[fx][f(f^n(x))],
\end{align*}
即 $\num{n+1}$。
\end{proof}

\begin{ex}
  我们来看看将 $\fn{Succ}$ 应用到~$\num{0}$ 上会发生什么，$\num{0}$ 即 $\lambd[fx][x]$。我们将项完整地写出来：
  \begin{align*}
    \fn{Succ}\, \num{0} & \ident (\lambd[a][\lambd[f][\lambd[x][(f (({a} f) x))]]])(\lambd[f][\lambd[x][x]])\\
    & \redone \lambd[f][\lambd[x][(f (({(\lambd[f][\lambd[x][x]])} f) x))]] \\
    & \redone \lambd[f][\lambd[x][(f ({(\lambd[x][x])} x))]] \\
    & \redone \lambd[f][\lambd[x][(f x)]] \ident \num{1}\\
  \end{align*}
\end{ex}

\begin{prob}
  项
  \begin{align*}
    \fn{Succ}' & \ident \lambd[{n}][\lambd[fx][{n} f (f x)]]
  \end{align*}
  $\lambd$-定义了后继函数。请说明其原理。
\end{prob}

\begin{prop}
  \ollabel{prop:add-ld}
  加法函数~$\Add$ 是 $\lambd$-可定义的。
\end{prop}

\begin{proof}
加法可由项
\begin{align*}
  \fn{Add} & \ident \lambd[{a}{b}][\lambd[fx][{a} f ({b} f x)]]
\intertext{或者，}
  \fn{Add}' & \ident \lambd[{a}{b}][{a}\, \fn{Succ} \, {b}].
\intertext{定义。第一个加法项的工作方式如下：$\fn{Add}$ 首先接受两个数字 ${a}$ 和 ${b}$ 作为自变元。结果是一个接受 $f$ 和 $x$ 并返回 $af(bfx)$ 的函数。如果 $a$ 和 $b$ 是丘奇数 $\num{n}$ 和 $\num{m}$，那么这归约为 $f^{n+m}(x)$，它与 $f^{n}(f^{m}(x))$ 等同。或者，一步步来看：}
  (\lambd[{a}{b}][\lambd[fx][{a} f ({b} f x)]])\num n\,\num m & \redone
  \lambd[fx][\num{n}\, f (\num {m}\, f x)] \\
  & \redone \lambd[fx][\num{n}\, f (f^m x)] \\
  & \redone \lambd[fx][f^n (f^m x)] \ident \num {n+m}.
\intertext{第二种加法表示 $\fn{Add'}$ 的工作方式不同：应用到两个丘奇数 $\num{n}$ 和 $\num{m}$ 上时，}
\fn{Add}' \num n \,\num m
& \redone \num{n}\, \fn{Succ}\, \num{m}.
\intertext{但 $\num{n} f x$ 总是归约为 $f^n(x)$。所以，}
  \num{n}\,  \fn{Succ}\, \num{m} & \red \fn{Succ}^n(\num{m}).
\end{align*}
由于 $\fn{Succ}$ $\lambd$-定义了后继函数，而对 $m$ 应用后继函数 $n$ 次得到 $n+m$，因此它最终归约为~$\num{n+m}$。
\end{proof}

\begin{prop}
  \ollabel{prop:mult-ld}
  乘法可由项
  \[
  \fn{Mult} \ident \lambd[ab][\lambd[fx][a (b f) x]]
  \]
  $\lambd$-定义。
\end{prop}

\begin{proof}
为了理解其工作原理，假设我们将 $\fn{Mult}$ 应用到丘奇数 $\num{n}$ 和 $\num{m}$ 上：$\fn{Mult} \, \num{n} \, \num{m}$ 归约为
  $\lambd[fx][\num{n}(\num{m}\, f)x]$。项 $\num{m} f$ 定义了一个将 $f$ 对其自变元应用 $m$ 次的函数。因此，$\num{n} (\num{m} f) x$ 将“将 $f$ 应用 $m$ 次”这个函数本身对 $x$ 应用 $n$ 次。换句话说，我们对 $x$ 应用 $f$ 共 $n\cdot m$ 次。但所得的范式正是丘奇数 $\num{nm}$。
\end{proof}

\begin{editorial}
实际上，我们可以通过 $\eta$-归约进一步简化这个项：
\[
  \fn{Mult} \ident \lambd[ab][\lambd[f][a (b f)]].
\]

但这需要我们首先解释 $\eta$-归约。
\end{editorial}

\begin{prob}
  乘法也可以由项
  \[
  \fn{Mult}' \ident \lambd[ab][a (\fn{Add}\, a) \num{0}].
\]
  $\lambd$-定义。
  请说明其原理。
\end{prob}

指数运算作为 $\lambd$-项的定义出奇地简单：
\[
  \fn{Exp} \ident \lambd[be][e b].
\]
第一个自变元 $b$ 是底数，第二个自变元~$e$ 是指数。直观上，根据我们对数字的编码，$e f$ 就是 $f^e$。如果你觉得这很难理解，我们仍然可以通过迭代乘法来定义指数运算：
\[
  \fn{Exp}' \ident \lambd[be][e (\fn{Mult}\, b) \num{1}].
\]

在丘奇数上定义前驱和减法并不像我们想象的那么简单：它需要用到序对的编码。

\end{document}